\documentclass[letterpaper,10pt]{article}
\usepackage[usenames,dvipsnames]{color}
\usepackage[default]{sourcesanspro}
\usepackage[hidelinks]{hyperref}
\usepackage[english]{babel}
\usepackage{amsmath}
\usepackage{csquotes}
\usepackage{fancyhdr}
\usepackage{enumitem}
\usepackage{latexsym}
\usepackage{marvosym}
\usepackage{soul}
\usepackage{tabularx}
\usepackage{titlesec}
\usepackage{verbatim}

% margin needs to be after empty
\usepackage[empty]{fullpage}
\usepackage[margin=0.5in]{geometry}

% Setting for Publications Page
\usepackage[style=ieee, sorting=ydnt, defernumbers=true, maxbibnames=4, minbibnames=4]{biblatex}
\input{glyphtounicode}
\setlength\bibitemsep{1\itemsep}
\addbibresource{publication.bib}

% ---------- PAGE SETUP ----------
\pagestyle{fancy}
\fancyhf{} % clear all header and footer fields
\fancyfoot{}
\setlength{\footskip}{4pt}

% remove the rule on the top
\renewcommand{\headrulewidth}{0pt}
\renewcommand{\footrulewidth}{0pt}
\raggedbottom
\raggedright

% Sections formatting
\titleformat{\section}{\vspace{-5pt}\scshape\raggedright\large}{}{0em}{}[\color{black}\titlerule \vspace{-4pt}]

% Ensure that generate pdf is machine readable/ATS parsable
\pdfgentounicode=1

%%%%%%  RESUME STARTS HERE  %%%%%%%%%%%%%%%%%%%%%%%%%%%%
\begin{document}
% ---------- HEADING ----------
\begin{flushleft}
  \textbf{\large Yixiang Gao},
  \textbf{\large PhD}, 
  \textit{\large Electrical and Computer Engineering}\\
  \vspace{2.5pt}
  \textbf{Contact:} 573-825-1296 $|$
  \href{mailto:yg5d6@umsystem.edu}{yg5d6@mail.missouri.edu}\\
\end{flushleft}

\begin{tabular*}{\textwidth}{@{\extracolsep{\fill}}l l}
  \href{https://www.linkedin.com/in/yixiang-gao-13a1408b/}
  {{\textbf{Linkedin:} https://tinyurl.com/yixiang-linkedin}} & 
  \href{https://scholar.google.com/citations?user=7104qXwAAAAJ&hl=en}
  {{\textbf{Google Scholar:} https://tinyurl.com/yixiang-googlescholar}}\\
  
  \href{https://g1y5x3.github.io}{{\textbf{Website:} https://g1y5x3.github.io}} &
  \href{https://github.com/g1y5x3/}{{\textbf{GitHub:} https://github.com/g1y5x3/}}
\end{tabular*} 

% ---------- PROFESSIONAL SUMMARY ----------
\section{Professional Summary}
\begin{itemize}[leftmargin=0.15in, label={}]
  \item{
    Machine learning and robotics researcher.
    Exhibits profound expertise in \textit{\textbf{Deep Learning}, \textbf{Computer Vision}, \textbf{Robotics} and \textbf{Pattern Recognition}}.
    Specialized in fostering interdisciplinary collaborations, notably with speech language pathology, and underground mining related applications. 
    Advocate for open source technologies.
    Seeking to leverage extensive AI/ML expertise to address complex real-world challenges and drive innovation in technology application.
  }
\end{itemize}

% ---------- EDUCATIONu ----------
\section{Education}
\begin{tabular*}{\textwidth}{@{\extracolsep{\fill}}l l r}
  University of Missouri - Columbia & 
  \textbf{Ph.D in Electrical and Computer Science} &
  \textit{2017 -- 2024} \\
  
  University of Missouri - Columbia &
  \textbf{BS in Electrical Engineering} &
  \textit{2014 -- 2017} \\
  
  University of Missouri - Columbia &
  \textbf{BS in Computer Engineering} &
  \textit{2014 -- 2017} \\
\end{tabular*}

% ---------- EXPERIENCE ----------
\section{Experience}

% Post Doctorial Fellow
\textbf{Post Doctorial Fellow} $|$
\emph{Missouri S\&T} $|$
\emph{Department of Mining \& Explosives Engineering} \hfill
\emph{Oct 2024 -- Current}
\begin{itemize}[leftmargin=0.15in, parsep=0pt] 
\item Robotics researcher for CDC NIOSH Initiative to empower miners for self-escape
  \begin{itemize}[leftmargin=0.15in, parsep=0pt] 
    \item Lead research on enhancing deep learning-based object detection algorithms such as YOLO in mine safety applications.
    \item Integrate advanced LiDAR and thermal camera on Boston Dynamics' Spot quadrupled robot with ROS2.
    \item Develop autonomous navigation solutions with Spot for underground miner exploration and miner safety monitoring.
  \end{itemize}
\end{itemize}

% Graduate Research Assistant
\textbf{Graduate Research Assistant} $|$
\emph{University of Missouri - Columbia} $|$
\emph{ViGIR Lab} \hfill
\emph{Aug 2017 -- Sep 2024}
\begin{itemize}[leftmargin=0.15in, parsep=0pt] 
  \item PhD Thesis: Confounded predictions in machine learning
  \begin{itemize}[leftmargin=0.15in, parsep=0pt] 
    \item Detect, quantify, and mitigate confounding factors in machine learning and deep learning models.
  \end{itemize}
  \item Student Investigator for NIH R01DC018026 \hfill
  \emph{Aug 2019 -- Sep 2024}
  \begin{itemize}[leftmargin=0.15in, parsep=0pt]
    \item Developed classification pipeline with transformers and CNN using Pytorch and Matlab for analyzing both sEMG and fMRI data. 
  \end{itemize}
  \item Student Investigator for NIH R15DC015335 \hfill
  \emph{Aug 2017 -- Aug 2019}
  \begin{itemize}[leftmargin=0.15in, parsep=0pt]
    \item Developed machine learning algorithms with SVM, K-means, and MLP using scikit-learn and Pytorch for analyzing sEMG data
  \end{itemize}
\end{itemize}

% Other projects
\textbf{Other projects}
\begin{itemize}[leftmargin=0.15in, parsep=0pt] 
  \item SpotMicro $|$ 
  \emph{University of Missouri - Columbia} $|$ 
  \emph{ViGIR Lab}
  \hfill\textit{2022 -- 2023}
  \begin{itemize}[leftmargin=0.15in, parsep=0pt]
    \item An open-source quadruped robot for the MU Robotics club. Operated by a Raspberry Pi 4 and integrated with ROS.
  \end{itemize}
  \item Object Detection and Pose Estimation using Embedded Devices $|$
   \emph{University of Missouri - Columbia} $|$ 
  \emph{ViGIR Lab}
  \hfill\textit{2019}
  \begin{itemize}[leftmargin=0.15in, parsep=0pt]
    \item Applied YOLO on a Raspberry Pi 3 with stereo-vision cameras for object detection and pose estimation.
  \end{itemize}
\end{itemize}

% Graduate Teaching Assistant
% \textbf{Graduate Research Assistant} $|$
% \emph{University of Missouri - Columbia} \hfill
% \emph{2017 -- 2019}
% \begin{itemize}[leftmargin=0.15in, parsep=0pt] 
%   \item Microprocessor Engineering
%   \item Software Design in C and C++
% \end{itemize}

% ---------- Publications ----------
\nocite{*}
\printbibliography[title={Publications}]

% ---------- TECHNICAL SKILLS ----------
\section{Technical Skills}
\begin{itemize}[leftmargin=0.15in, parsep=0pt, label={}]
  \item \textbf{Programming Languages: }{Python, C/C++/CUDA, Matlab, HTML}
  \item \textbf{Libraries: }{pytorch, tinygrad, scikit-learn, transformers, ultralytics, mmdet, ROS 2, Gazebo, ONNX}
  \item \textbf{Experienced Models: }{SVM, GA, MLP, CNN, RNN, transformers, ViT, DETR, YOLOs, Llama}
  \item \textbf{MLOps and API Tools: }{Git, Docker, Kubernetes, Github Actions for CI/CD}
  \item \textbf{Operating Systems: }{Linux(preferred), Windows}
  \item \textbf{Languages: }{English, Chinese}
\end{itemize}

% ---------- Other activites ----------
\section{Other activities}
\begin{itemize}[leftmargin=0.15in, label={}]
  \item{
  \textbf{Academic Paper Reviews: }{EMBC2019/2020/2021, CEC2023, Applied Sciences, TAI} \\
  \textbf{Open-Source Project Contributions: }{tinygrad, huggingface transformers} \\
  }
\end{itemize}

%----------HEADING---------------------
%\section{Awards and Certifications}
%    \resumeSubHeadingListStart
%      \resumeProjectHeading
%          {\titleItem{Star Award} \emph{ $|$ Company 1}}{September 2018}
%      \resumeProjectHeading
%          {\titleItem{Kudos Award} \emph{ $|$ Company 2}}{May 2017}
%      \resumeProjectHeading
%          {\titleItem{IBM Full Stack Software Developer} \emph{$|$ IBM}}{November 2023}
%          \resumeProjectHeading
%          {\titleItem{CompTia Sec+} \emph{$|$ CompTia}}{December 2023}
%    \resumeSubHeadingListEnd
%-------------------------------------------

\end{document}
